\begin{pastexam}[label={px:nsysu-114-osds-01}]{中山資工所甲組 114 學年度 第 1 題}
（追蹤）

What is printed by each of the following C program?
\begin{enumerate}[label=(\alph*)]
\item (6\%)
\begin{minted}[fontsize=\small]{c}
int a=2, b=13, c;
for (c=1; b>0; b = b>>1) {
   if (b%2 == 1) c = c*a;
   a = a*a;
}
printf("%d \n", c);
\end{minted}
\item (9\%)
\begin{minted}[fontsize=\small]{c}
int a[64]={0,1,2,3,4,5,6,7,8,9,10,11,12,13,14,15};
int c=1;
for (int i=1; i<=3; i++) {
     for(int j=0; j<=15; j++)
         a[j] = a[j]+a[j+c];
     c *= 2;
}
printf("%d %d %d \n", a[0], a[5], a[14]);
\end{minted}
\end{enumerate}
\end{pastexam}

\begin{pastsoln}
答案：(a) 印出 \texttt{8192}；(b) 印出 \texttt{28 68 29}。

(a) 是資結書第~17.1.5~節的快速冪（迴圈版），(b) 是陣列上的倍增累加（資結書第~4~章的陣列操作）；
本題放在第~4~章。

\textbf{(a)} 迴圈每一輪看 $b$ 的最低位元：是 $1$ 就把目前的 $a$ 乘進 $c$；然後 $a$ 自乘、$b$ 右移一位。 所以第 $k$ 輪
（從 $0$ 數）的 $a$ 是 $2^{2^k}$，$c$ 乘進的是 $b$ 的二進位中為 $1$ 的那些位所對應的 $2^{2^k}$。 $13 = 1101_2$：
\begin{center}
\begin{tabular}{ccccc}
\toprule
輪 & $b$（二進位） & $b \,\%\, 2$ & $c$ & 之後的 $a$ \\
\midrule
0 & 1101 & 1 & $1 \times 2 = 2$ & $4$ \\
1 & 110 & 0 & $2$ & $16$ \\
2 & 11 & 1 & $2 \times 16 = 32$ & $256$ \\
3 & 1 & 1 & $32 \times 256 = 8192$ & $65536$ \\
\bottomrule
\end{tabular}
\end{center}
之後 $b = 0$，迴圈結束。 $c = 2^{1} \cdot 2^{4} \cdot 2^{8} = 2^{13} = 8192$，程式算的就是 $a^b$，迴圈只跑
$\lfloor \lg b \rfloor + 1$ 次。 最後一輪的 $a = 65536$ 沒有用到，也沒有溢位（\texttt{int} 至少 $32$ 位元時）。
輸出 \texttt{8192}（後面跟一個空白和換行）。

\textbf{(b)} 陣列只有 $a[0..15] = 0..15$，其餘 $a[16..63]$ 都是 $0$，而且內層迴圈只改 $a[0..15]$，所以
$a[16..19]$（會被讀到的最大索引是 $15 + 4$）一直是 $0$。

\textbf{關鍵觀察。} 內層迴圈 $j$ 由小到大，$a[j]$ 讀的是 $a[j + c]$，$j + c > j$，這一格在本輪還沒被改過，所以每一輪
等於用\textbf{上一輪}的值同時做 $a[j] \gets a[j] + a[j + c]$。 令原本的值 $x_k = k$（$k \le 15$）、$x_k = 0$（$k \ge 16$）：
\begin{itemize}
\item 第 1 輪（$c = 1$）後：$a[j] = x_j + x_{j+1}$，連續 $2$ 個的和；
\item 第 2 輪（$c = 2$）後：$a[j] = (x_j + x_{j+1}) + (x_{j+2} + x_{j+3})$，連續 $4$ 個；
\item 第 3 輪（$c = 4$）後：連續 $8$ 個，$a[j] = \sum_{k=j}^{j+7} x_k$。
\end{itemize}
每一輪把區間長度加倍。 逐輪的 $a[0..15]$：
\begin{center}
\small
\setlength{\tabcolsep}{3pt}
\begin{tabular}{l*{16}{r}}
\toprule
$j$ & 0 & 1 & 2 & 3 & 4 & 5 & 6 & 7 & 8 & 9 & 10 & 11 & 12 & 13 & 14 & 15 \\
\midrule
第 1 輪 & 1 & 3 & 5 & 7 & 9 & 11 & 13 & 15 & 17 & 19 & 21 & 23 & 25 & 27 & 29 & 15 \\
第 2 輪 & 6 & 10 & 14 & 18 & 22 & 26 & 30 & 34 & 38 & 42 & 46 & 50 & 54 & 42 & 29 & 15 \\
第 3 輪 & \textbf{28} & 36 & 44 & 52 & 60 & \textbf{68} & 76 & 84 & 92 & 84 & 75 & 65 & 54 & 42 & \textbf{29} & 15 \\
\bottomrule
\end{tabular}
\end{center}
用公式核對三個要印的值：
\[
a[0] = 0 + 1 + \cdots + 7 = 28,\quad a[5] = 5 + 6 + \cdots + 12 = 68,\quad a[14] = 14 + 15 = 29
\]
（$a[14]$ 的區間 $14..21$ 裡只有 $14$、$15$ 不是 $0$）。 輸出 \texttt{28 68 29}。
\end{pastsoln}
